\documentclass[letterpaper]{article}
\usepackage{amsmath,amssymb}
\usepackage[letterpaper,margin=1in]{geometry}
\setlength{\parindent}{0pt}

\begin{document}

\textbf{MATH 1564 --- Notes 09/24}

\bigskip

\[
\frac{F_k}{r_+^k}\longrightarrow c_+ \qquad \text{(homework)}
\]

\[
F_k\text{ grows like }c_+r_+^k
=c_+\left(\frac{1+\sqrt5}{2}\right)^k.
\]

\bigskip

\textbf{Theorem.} Let
\[
\operatorname{Fib}=\left\{(x_k)_{k\geq0}:x_{k+2}=x_{k+1}+x_k\right\}.
\]
\begin{enumerate}
\item $\operatorname{Fib}$ is a subspace of $S$.
\item The map $T:\operatorname{Fib}\to\mathbb{R}^2$, given by
\[
T\bigl((x_k)\bigr)=\begin{bmatrix}x_0\\x_1\end{bmatrix},
\]
is linear and invertible.
\item $\{(r_+^k),(r_-^k)\}$ is a basis in $\operatorname{Fib}$, where
\[
r_\pm=\frac{1\pm\sqrt5}{2}.
\]
\item Any vector $(x_k)$ in $\operatorname{Fib}$ can be written as
\[
c_+(r_+^k)+c_-(r_-^k),
\]
where
\[
\begin{bmatrix}c_+\\c_-\end{bmatrix}
\]
are coordinates of $(x_k)$ in the basis $\{(r_+^k),(r_-^k)\}$.
\end{enumerate}

\bigskip

\textbf{Generalization to 2nd-order recursions}

\medskip

Let $(x_k)$ be a sequence in $S$ such that
\[
x_{k+2}+a_1x_{k+1}+a_0x_k=0\quad \text{for all }k,\qquad a_0\ne0.
\]

\textbf{Theorem.} Such sequences form a 2-dimensional subspace $H$ of $S$, the vector space of all sequences.

\textbf{Proof.} Let $T:H\to\mathbb{R}^2$ be
\[
T\bigl((x_k)_{k\geq0}\bigr)=\begin{bmatrix}x_0\\x_1\end{bmatrix}.
\]
Then $T$ is linear and invertible.

\bigskip

\textbf{How to find solutions.} Consider
\[
x_{k+2}+a_1x_{k+1}+a_0x_k=0,\qquad a_0\ne0.
\]
Search for solutions of the form $(r^k)_{k\geq0}$ for some $r\ne0$. Then
\[
r^{k+2}+a_1r^{k+1}+a_0r^k=0.
\]
Dividing by $r^k$ gives
\[
r^2+a_1r+a_0=0,
\]
so
\[
r=\frac{-a_1\pm\sqrt{a_1^2-4a_0}}{2}.
\]

\newpage

\textbf{Case 1:} $a_1^2>4a_0$.

\medskip

The two nonzero roots $r_+$ and $r_-$ are two independent vectors in the 2-dimensional space $H$. They form a basis, so any vector in $H$ is
\[
c_+(r_+^k)+c_-(r_-^k).
\]

\textbf{Example.} $a_1=-1=a_0$: Fibonacci.

\bigskip

\textbf{Case 2:} $a_1^2=4a_0$.

\[
r=-\frac{a_1}{2}\ne0.
\]
Then $(r^k)_{k\geq0}=\left(\left(-\frac{a_1}{2}\right)^k\right)_{k\geq0}$ solves the recursion. In this case, another vector is needed to form a basis in $H$.

Try
\[
(x_k)_{k\geq0}=(kr^k)_{k\geq0},
\qquad r=-\frac{a_1}{2}.
\]

\textbf{Substitution.}
\begin{align*}
&(k+2)r^{k+2}+a_1(k+1)r^{k+1}+a_0kr^k\\
&=r^k\left[(k+2)r^2+a_1(k+1)r+a_0k\right]\\
&=r^k\left[(k+2)\frac{a_1^2}{4}-\frac{a_1^2}{2}(k+1)+\frac{a_1^2}{4}k\right]\\
&=r^k\frac{a_1^2}{4}\left[k+2-2(k+1)+k\right]=0.
\end{align*}
Here $r=-a_1/2$ and $a_0=a_1^2/4$.

\bigskip

\textbf{Theorem.} If $r^2+a_1r+a_0=0$, $a_0\ne0$, has a repeated root
\[
r=-\frac{a_1}{2}\ne0,
\]
then $\{(r^k),(kr^k)\}$ is a basis in the solution space $H$.

\textbf{Example.}
\[
x_{k+2}+2x_{k+1}+x_k=0,
\qquad r^2+2r+1=0,
\qquad r=-1.
\]
Thus
\[
\{(-1)^k,k(-1)^k\}
\]
is a basis in $H$.

\newpage

\textbf{Case 3:} $a_1^2<4a_0$.

The equation $r^2+a_1r+a_0=0$ has two complex roots:
\[
r_\pm=\frac{-a_1\pm\sqrt{a_1^2-4a_0}}{2}
=\frac{-a_1\pm i\sqrt{4a_0-a_1^2}}{2}.
\]
Then $(r_+^k)$ and $(r_-^k)$ solve
\[
x_{k+2}+a_1x_{k+1}+a_0x_k=0.
\]

\textbf{Example.} If $a_1=0$ and $a_0=1$, then
\[
r_\pm=\pm i,
\]
and $(i^k),((-i)^k)$ solve the recursion
\[
x_{k+2}+x_k=0.
\]
Also,
\[
i^k=(i,-1,-i,1,i,i^2,i^3,i^4,\ldots).
\]

\bigskip

If $r_+^k,r_-^k$ are nonreal solutions of
\[
x_{k+2}+a_1x_{k+1}+a_0x_k=0,
\]
then
\[
\frac{r_+^k+r_-^k}{2},
\qquad
\frac{r_+^k-r_-^k}{2i}
\]
are real solutions.

\textbf{Proof.} Since $r_-=\overline{r_+}$,
\[
r_-^k=\overline{r_+^k}.
\]
Writing $r_+^k=a+ib$ gives
\[
\frac{r_+^k+r_-^k}{2}
=\frac{a+ib+a-ib}{2}=a,
\]
and
\[
\frac{r_+^k-r_-^k}{2i}
=\frac{a+ib-(a-ib)}{2i}=b.
\]

\bigskip

\textbf{Theorem.} The two real solutions
\[
\left(\frac{r_+^k+r_-^k}{2}\right),
\qquad
\left(\frac{r_+^k-r_-^k}{2i}\right)
\]
are independent in $S$.

\textbf{Proof.} Let $u=(r_+^k)$. Then
\[
(r_-^k)=\overline{(r_+^k)}=\overline{u}.
\]
Suppose there is a real number $c$ with
\[
\frac{u+\overline{u}}{2}=c\frac{u-\overline{u}}{2i}.
\]
Then
\[
iu+i\overline{u}=cu-c\overline{u},
\]
so
\[
(c+i)\overline{u}=(c-i)u=(c-i)(r_+^k).
\]
Thus
\[
\left(\frac{r_+}{r_-}\right)^k=\frac{c+i}{c-i}
\]
is constant, which is impossible.

\newpage

\textbf{Theorem.} If $r_+,r_-$ are nonreal distinct roots of
\[
r^2+a_1r+a_0=0,
\]
then
\[
\left\{
\frac{r_+^k+r_-^k}{2},
\frac{r_+^k-r_-^k}{2i}
\right\}
\]
is a basis in the solution space of
\[
x_{k+2}+a_1x_{k+1}+a_0x_k=0.
\]

\bigskip

\textbf{Generalization}

Let $A$ be an $n\times n$ real matrix. Let $(X_k)_{k\geq0}$ be a sequence of vectors in $\mathbb{R}^n$.

Let $S^n$ be the vector space of sequences of vectors in $\mathbb{R}^n$:
\[
(X_1,AX_1,A^2X_1,A^3X_1,\ldots).
\]

\textbf{Recursion.}
\[
X_{k+1}=AX_k.
\]

\textbf{Example.} Write
\[
x_{k+2}+a_1x_{k+1}+a_0x_k=0
\]
in matrix form. Let
\[
A=\begin{bmatrix}0&1\\-a_0&-a_1\end{bmatrix},
\qquad
X_k=\begin{bmatrix}x_k\\x_{k+1}\end{bmatrix},
\qquad
X_{k+1}=\begin{bmatrix}x_{k+1}\\x_{k+2}\end{bmatrix}.
\]
Then
\[
X_{k+1}=AX_k
=\begin{bmatrix}0&1\\-a_0&-a_1\end{bmatrix}
\begin{bmatrix}x_k\\x_{k+1}\end{bmatrix}
=\begin{bmatrix}x_{k+1}\\-a_0x_k-a_1x_{k+1}\end{bmatrix}.
\]

\textbf{Theorem.} Let $A$ be an $n\times n$ real matrix. The solution space of
\[
X_{k+1}=AX_k
\]
has dimension $n$.

\textbf{Proof.}
\[
X_{k+1}=AX_k=AAX_{k-1}=\cdots=A^kX_1.
\]
Let $T:S^n\to\mathbb{R}^n$ be the map sending the sequence $(X_k)_{k\geq1}$ to $X_1$. The map is invertible, since the linear map $T$ is a linear isomorphism.

\newpage

\textbf{Notation.}
\[
z=a+ib,\qquad a,b\text{ real}.
\]
One can do linear algebra with matrices whose entries are complex numbers:
\[
\mathbb{C}^n=\{\text{all }n\times1\text{ matrices with complex entries}\}.
\]
One can talk about vector spaces with complex scalars.

\textbf{Examples.} Polynomials with complex coefficients; sequences $(x_k)$ where each $x_k$ is a complex number.

\bigskip

\textbf{Theorem.} Complex conjugation commutes with addition and multiplication of complex numbers.

\textbf{Corollary.} Complex conjugation commutes with addition and multiplication of matrices.

\textbf{Example.}
\[
\overline{\begin{bmatrix}i&2+i\\-1&-1-i\end{bmatrix}}
=\begin{bmatrix}-i&2-i\\-1&-1+i\end{bmatrix}.
\]

\textbf{Proof.} The $ij$ entry of $\overline{A+B}$ is the conjugate of the $ij$ entry of $A+B$, so
\[
\overline{\text{$ij$ entry of }A+\text{$ij$ entry of }B}
=\text{$ij$ entry of }\overline{A}+\text{$ij$ entry of }\overline{B}.
\]
For multiplication, the $ij$ entry of $\overline{A}\,\overline{B}$ is
\[
\begin{bmatrix}\overline{a_{i1}}&\cdots&\overline{a_{in}}\end{bmatrix}
\begin{bmatrix}\overline{b_{1j}}\\\vdots\\\overline{b_{nj}}\end{bmatrix}
=\overline{a_{i1}}\,\overline{b_{1j}}+\cdots+\overline{a_{in}}\,\overline{b_{nj}}
=\overline{a_{i1}b_{1j}+\cdots+a_{in}b_{nj}},
\]
which is the $ij$ entry of $\overline{AB}$.

\newpage

\textbf{Example: $n$th-order linear recursion}

Let $(x_k)$ be a sequence of real numbers that solves
\[
x_{k+n}+a_{n-1}x_{k+n-1}+\cdots+a_0x_k=0,
\qquad a_0\ne0. \tag{*}
\]
The space of such sequences has dimension $n$.

Let
\[
X_k=\begin{bmatrix}x_k\\x_{k+1}\\\vdots\\x_{k+n-1}\end{bmatrix},
\qquad
X_{k+1}=\begin{bmatrix}x_{k+1}\\x_{k+2}\\\vdots\\x_{k+n}\end{bmatrix},
\]
and
\[
A=
\begin{bmatrix}
0&1&0&\cdots&0\\
\vdots&\ddots&\ddots&\ddots&\vdots\\
\vdots&&\ddots&\ddots&0\\
0&\cdots&0&0&1\\
-a_0&-a_1&\cdots&-a_{n-2}&-a_{n-1}
\end{bmatrix}.
\]
Then $X_{k+1}=AX_k$.

\bigskip

\textbf{Theorem.} The map $(X_k)\mapsto(x_k)$ is a linear isomorphism, so the solution space of $(*)$ has dimension $n$.

We assume $a_0\ne0$ because \ldots

\bigskip

\textbf{Example.} What is the dimension of the solution space of
\[
x_{k+4}-x_{k+2}+x_{k+1}=0\,?
\]
The index $k$ appears in the recursion only through $x_{k+1},x_{k+2},x_{k+4}$, so $x_0$ never enters it. Shifting the index (working with $(x_k)_{k\geq1}$) gives
\[
x_{k+3}-x_{k+1}+x_k=0,
\]
a 3rd-order recursion with $a_0=1\ne0$, so the theorem applies:
\[
\dim=3.
\]

\bigskip

\textbf{Note.} We need to learn more linear algebra to solve $X_{k+1}=AX_k$. We need to compute $A^k$ for all $k$ at once.

\end{document}
